% ===== §1 =====
\section{Introduction: An Operation Mistaken for a Law}
\label{sec:intro}

\noindent
A question can be so widely presupposed that it is never asked, and the presupposition of this paper is one of these. In the ordinary speech of the present, a certain family of words recurs with the confidence of the self-evident. One speaks of the \emph{optimal} choice, of the \emph{efficient} arrangement, of whether a thing was \emph{worth it}, of getting the \emph{most} out of a day, a career, a marriage. The words are not idle. They carry a picture, and the picture has hardened, over the century in which it became the native tongue of economics, management, and much of the science of decision, into something close to a metaphysics. On this picture there is, for any situation in which one might act, a quantity to be made as large as possible; the situations of a life are so many landscapes, and to live well is to climb them; and the good, wherever it is found, is whatever sits at the summit. The picture is so pervasive that its central operation, the taking of a maximum, has ceased to look like one operation among others with a domain of its own and has come to look like the form of rationality as such, and beyond that, dimly, like the form of things. This paper is about that hardening, and its guiding claim is a denial: that optimization is not a law of the universe, and that morality and the manner of being of intimate relation are not, in their structure, optimization problems at all.

The denial is easy to state and easy to mistake. It is not the pious complaint that one \emph{ought not} to optimize love, that to do so is cold or unromantic. Such a complaint concedes the very thing in dispute, for it grants that an optimization problem is there to be solved and merely asks that we refrain from solving it, out of tact. The claim of this paper is stronger and colder than that. It is that in intimate relation there is, in the strict sense, no well-defined optimization problem to refrain from: that the operation of maximization, applied here, does not have a hard solution one declines to compute but has no legitimate solution to compute at all, because the conditions under which such an operation is defined are not, in this domain, satisfied. The economist who hears the complaint has a ready reply: that love resists optimization only because our utility function is not yet fully written, and that the remedy is a richer function, one with intimacy and meaning and the long horizon folded in. The reply is powerful and must be met on its own ground. The burden of the paper's central part is to show that it fails, that the trouble is not a function written too poorly but the impossibility, for reasons given in their place, of there being a function to write.

\subsection*{The origin of the question, and its stakes}

The question that organizes the paper arrived from without, in the course of an exchange on the political economy of value and the theory of generative justice, and it arrived in a form so simple that it took some time to feel its weight: whether optimization really is a law of the universe. The interlocutor's own framework had long resisted the reduction of heterogeneous value to a single settling scale, and the question was in that sense already implicit in it; but stated baldly it did a particular work. It moved optimization from the place of an obvious and unexamined procedure to the place of a metaphysical thesis, which is to say a thesis that might be false. Once it is seen as a thesis, one can ask what it asserts and whether it is true, and the asking is the whole of this paper. The stakes are not confined to intimacy. If the argument holds, then a picture that governs a great deal of contemporary reasoning about how to live is shown to rest, at least in one central region, on a category mistake; and intimate relation, which the series has taken throughout as the minimal and most exposed unit of relational being, is the region in which the mistake is most legible, because it is there that the goods which resist the picture are least disguised.

\subsection*{Two movements, kept deliberately apart}

The paper's argument falls into two movements which it is essential not to confuse, and much of the force of the whole depends on holding them apart. The first movement is a proof, and it is internal and formal. It asks whether the intimate relation furnishes a well-defined optimization problem, and it answers, in four steps, that it does not; and it does this, deliberately, \emph{without any appeal to value}. It does not say that optimizing is bad. It says only that the operation is here undefined: that there is nothing of the required shape for the operation to act upon. One may read the whole of the second part of this paper as an atheist would read a proof that a certain equation has no solution over the reals, with no moral coloring at all, and the argument should be able to compel a reader who feels no particular reverence for love and merely wishes his formalisms to be well posed.

The second movement is a criticism, and it is external and normative. It takes the result of the first as given, that here there is no optimization problem, and it asks what follows when the picture is imposed regardless, when a domain that furnishes no well-defined maximization is treated as though it did. Here, and only here, do the value-laden claims enter: that the imposition damages what it touches, converting a generative circulation into an extractable and exhaustible quantity; that it is a symptom of a particular historical condition rather than a timeless truth; and that it mistakes the kind of thing morality is. The order matters. The proof clears the ground by showing that the frame does not fit; the criticism then examines the cost of forcing it on. To run them together, to argue at once that optimization is undefined and that it is wicked, would weaken both, for it would let the reader suspect that the formal claim is only the moral distaste in disguise. It is not, and the separation is the guarantee.

\subsection*{Three lines, and the shape of the recovery}

Three lines run through the argument, as they have run through the series, and naming them here lets their recurrence read as design. The first is \emph{psychoanalytic}: that among the goods any complete objective would have to weigh is the drive, and that the drive, arising at a coupling of the Real and the Symbolic, is precisely what cannot be fully written down, so that the non-structurability of value is not a limitation of our present formalisms but a feature of the thing \citep{lacan1977xi}. The second is the line of \emph{relational ontology and the vacancy of the big Other}, developed across the series and stated most sharply in the papers on the non-binding vow and on relational estrangement: that value is not a property held inside a subject but is generated within a relation and exists only there, so that there is no position outside the relation, no omniscient third, from which its value could be read off and fixed \citep{paperxiii}. This line does the heaviest work, for it supplies the reason the objective must be endogenous, and it ties the whole to the series' recurring insistence that there is no metalanguage and no big Other to occupy the God's-eye view that quantification silently requires. The third is the line of \emph{generative justice and political economy}, developed in exchange with Eglash: that value which should circulate and return to those who generate it can instead be settled and drawn off, and that the imposition of an optimizing frame is one of the forms this extraction takes \citep{eglash2016}.

These lines converge on a recovery, for the paper does not end in denial. Having argued that rationality is not optimization, it owes an account of what relational rationality then is, and it offers one: a rationality of response, of attunement, and of the sustaining of a generation, which has its own severity and is not the severity of the maximum. But the recovery is held, in the manner the series demands, short of the temptation that would undo it. It would be easy, and it would be a defeat, to answer the question ``what should one optimize, if not utility?'' with ``love, then, or meaning, or generativity,'' for that answer keeps the frame and merely repaints the summit. The paper refuses it. Its positive claim is that relational rationality is plural and admits no synthesis into a single measure, and that the refusal to name one master objective is not an evasion but the content of the view, enacted, as the series enacts its claims, in the polyphonic form of its own close.
